A centered interaction table lives in a linear space.  A joint probability table occupies a
constrained subset: its entries remain nonnegative, sum to one, and often preserve two prescribed
marginals.  These constraints form a transport polytope.  The two geometries describe an
infinitesimal interaction coordinate and a finite probabilistic model.

Begin with independent marginals $p_i(a)p_j(b)$ and tilt them by a pair score.  At weak coupling,
double centering removes directions that only change one marginal, leaving the tangent interaction
inside the zero-marginal subspace.  At finite coupling, however, linear projection is insufficient:
exponentiation bends the geometry and generally changes both marginals.

Exact marginal preservation requires multiplicative row and column factors.  Matrix scaling finds
those factors and returns the unique information projection onto the transport polytope under the
usual positivity assumptions.  The nonlinear construction and the linear gauge are complementary
resolutions: the gauge-fixed operator is the local tangent coordinate through which the
probability-exact family passes.

This relation matters whenever a pair table is given probabilistic language.  Tangent-space norms
and finite-strength mutual information are different functionals.  A centered logit table reaches
a joint law with desired marginals through matrix scaling, and separately fitted conditionals reach
one global distribution only when their compatibility equations hold.  Probability geometry
supplies these finite constraints around the local operator coordinate.

The chapter develops that distinction without abandoning the operator view.  Exact pair laws,
information projections, and tangent operators describe different resolutions of the same local
departure from independence.  Keeping the resolution explicit prevents a convenient linear score
from being mistaken for a fully normalized physical model.
